


\section{Requêtes avec classement}

\begin{frame}{Classement par pertinence}

Les requêtes Booléennes classent les documents en deux catégories\,: ceux qui satisfont la requête, et les autres. C'est 1, ou c'est 0.

\vfill

Le \alert{classement par pertinence} (\textit{ranked search}) trie un résultat en fonction d'un ``poids'' (\textit{weight}) 
mesurant le degré de pertinence d'un document pour une recherche.


\vfill

Plusieurs approches complémentaires pour évaluer la pertinence d'un document $d$ pour une recherche $q$\,:

\begin{itemize}
  \item par \alert{similarité} entre $d$ et $q$\,;
  \item par \alert{l'importance} de $d$, évaluée par exemple par
       sa position dans une collection structurée  (e.g., graphe, cf. PageRank)\,;
  \item en prenant en compte  l'utilisateur (\alert{profil}, \alert{boucles de pertinence}, etc.). Cf. l'analyse des \alert{clics}\,!
\end{itemize}
\end{frame}

\begin{frame}{Recherche par similarité}


La requête $q$ \alert{et} le document $d$ sont placés 
     dans un même \alert{Espace métrique}, doté d'une fonction de distance $m_E$. 

\vfill

Une fonction $f$ produit un \alert{descripteur} (le plus souvent sous la forme d'un vecteur
appelé \textit{features vector})
à partir d'un document.
     
\vfill

La \alert{similarité}, ou \alert{score}, est l'inverse de la distance.

\[sim(q,d) = \frac{1}{m_E(f(q),f(d))}\]


\begin{block}{Langage de requêtes ultra-simplifié}
La requête est exprimée par un ensemble de mots-clés, et interprétée comme un "document" simplifié.
\end{block}

\end{frame}


\begin{frame}{Quelques caractéristiques}

Avec l'approche par similarité, le résultat (c.à.d. les documents qui ont un score non nul) 
est potentiellement \alert{très grand}.

\vfill

Il est impératif de \alert{classer} le résultat par ordre de score croissant, et
de présenter les $k$ prtemiers à l'utilisateur (typ., $k \simeq 10-20$).

\vfill

Souvenez-vous: rappel et précision 

\begin{itemize}
  \item La \alert{précision} est la fraction du résultat qui est vraiment pertinente.
  
  $$precision =\frac{|relevant| \cap |retrieved|}{|retrieved|}$$

\item Le \alert{rappel} est la fraction des documents pertinents présente dans le résultat.
$$recall =\frac{|relevant| \cap |retrieved|}{|relevant|}$$
\end{itemize}
\end{frame}




\begin{frame}{Vocabulaire usuel}

On place \alert{documents} et \alert{requête} dans un \alert{espace métrique}, souvent un \alert{espace vectoriel}.

\vfill

L'essentiel: une fonction de \alert{similarité}, définie à partir d'une distance, ou directement.

\vfill

Les éléments de cet espace sont appelés \alert{descripteurs} ou \alert{vecteurs de caractéristiques}
(\textit{features vector}).

\vfill

Le \alert{score} $sim(q,d)$ mesure la pertinence d'un document $d$ par rapport à une requête (besoin) $q$.

\vfill

Le résultat est \alert{classé} sur le score; les $k$ premiers documents sont présentés.

\end{frame}

\section{Première approche, recherche plein texte}



\begin{frame}{Soyons concrets: première approche pour la recherche plein texte}

Supposons connu l'ensemble de tous les termes $V=\{t_1, t_2, \cdots, t_n\}$
 de tous les termes utilisables pour la rédaction d'un
document.

\vfill

Exemple: $V$ = \{"papa", "maman", "gateau", "chocolat", "haut", "bas"\}

\vfill

On définit $E=\{0, 1\}^{n}$ comme l'espace de tous les vecteurs
constitués de $n$ coordonnées valant soit 0, soit 1. Ce sont nos descripteurs.

\vfill

Exemple: vecteurs  constitués de 6 coordonnées valant soit 0, soit 1. 

\vfill

Fonction $f$ associant  un document $d$ à son descripteur $v = f(d)$.

\[
      v[i] = \left\{
                \begin{array}{ll}
                  1 \text{ si $d$ contient le terme $t_i$}\\
                  0 \text{ sinon}
                \end{array}
              \right.
\]

\begin{block}{Jusque là, rien de nouveau}
C'est la représentation déjà vue pour les matrices d'incidences.
\end{block}
\end{frame}

\begin{frame}{Exemple(s)}

Rappel: on a $V$ = \{"papa", "maman", "gateau", "chocolat", "haut", "bas"\}

\vfill


Soit le document $d_{maman}$=\texttt{maman est en haut, qui fait du gateau}

\vfill

alors $f(d_{maman}) = [0, 1, 1, 0, 1, 0]$

\begin{exoblock}
À vous de jouer: $d_{papa}$= \texttt{papa est en bas, qui fait du chocolat}
\end{exoblock}
 
\vfill

Deux remarques:

\begin{itemize}
  \item On ignore certains mots (les "mots inutiles" ou \alert{stop words})
  \item L'ordre des mots dans le document est \alert{ignoré} (approche "\textit{bag of words}")
\end{itemize}

\end{frame}

\begin{frame}{La distance}

Dans un espace vectoriel, on peut penser à prendre la \alert{distance Euclidienne}. Si
$v_1$ et  $v_2$  sont deux vecteurs:

   $$E(v_1, v_2) = \sqrt{(v_1^1 - v^1_2)^2 + (v^2_1 - v^2_2)^2 + \cdots + (v_1^n - v^n_2)^2}$$

\vfill

Et la similarité est l'inverse de la distance

\[
      sim(v_1, v_2) = \left\{
                \begin{array}{ll}
                  \infty \text{ si $v^i_1=v^i_2$ pour tout $i$ }\\
                  \frac{1}{E(v_1, v_2)}  \text{ sinon}
                \end{array}
              \right.
\]

\vfill

Exemple, pour $q$ = "\texttt{maman haut chocolat}", $v_q = [0, 1, 0, 1, 1, 0]$

  $$sim (v_q, d_{maman}) = \frac{1}{\sqrt{2}}$$
  
\begin{exoblock}
À vous de calculer  $sim (v_q, d_{papa})$
\end{exoblock}
\end{frame}

\begin{frame}{Quelques points à retenir}

\begin{block}{Important}
Une différence concrète très sensible (illustrée ci-dessus) avec 
les requêtes Booléennes est qu’il n’est pas nécessaire qu’un document contienne tous les termes de la requête pour que son score soit différent de 0.
\end{block}

\begin{itemize}
   \item   "chocolat", un des mots-clés de $q$, 
   n’apparaît pas dans le document $d_{maman}$, malgré tout classé en tête;
\end{itemize}

\vfill

Approche présenté ici: simple mais \alert{nombreux inconvénients} (exercices).  

\vfill

Ignore l'impact de :  la taille des documents,
la taille du vocabulaire, le nombre d'occurrences d'un terme dans un document
et la rareté de ce terme dans la collection.

\end{frame}
